\section*{الفصل \ref{ch:b2:energy-geology} --- جيولوجيا الطاقة}
% numbers from the Book 4 second-writer scratch calculator calc/ch30.py

\begin{solution}{exo:b2:energy-geology:1}
(أ) 1.58: النوع~I؛ و0.95: النوع~III؛ و1.24: النوع~II؛ و1.12: نوع مختلط II/III. والكيروجين من النوع~III،
عند 0.95، هو الذي يعطي الغاز أساسًا. (ب) \mbox{\textdir TLT H/C $= 150/100 = 1.50$}، و\mbox{\textdir TLT O/C $= 8/100 = 0.08$}. (ج) لكلّ \qty{100}{g}:
$77.4/12.011 = \qty{6.444}{mol}$ من الكربون، و$9.8/1.008 = \qty{9.722}{mol}$ من الهيدروجين، و$8.2/15.999 =
\qty{0.513}{mol}$ من الأكسجين؛ فتكون \mbox{\textdir TLT H/C $= 1.51$} و\mbox{\textdir TLT O/C $= 0.0795$}، وهذا يوافق الصيغة.
\end{solution}

\begin{solution}{exo:b2:energy-geology:2}
(أ) 0.45: غير ناضجة، دون البداية عند \qty{0.60}{\percent}؛ و0.68: تولّد النفط؛ و1.1: لا تزال في
نافذة النفط، قرب نهايتها؛ و1.6: تجاوزت 1.20--\qty{1.30}{\percent}، فيتفكّك النفط إلى ميثان
ويعطي الصخر غازًا؛ و2.5: غاز فقط، ويستمرّ التوليد حتى نحو \qty{4.8}{\percent}. (ب) لا: فتوليد
الميثان من النوع~III يبدأ قرب \qty{0.73}{\percent}. (ج) $F = (\ln R_o + 1.6)/3.7$: $(-0.511 +
1.6)/3.7 = 0.294$ و$(0.262 + 1.6)/3.7 = 0.503$.
\end{solution}

\begin{solution}{exo:b2:energy-geology:3}
(أ) في المكمن، $6 \times 10^6 \times 25 \times 0.18 \times 0.70 = \qty{1.89e7}{m^3}$؛
و$\mathrm{STOIIP} = 1.89 \times 10^7/1.25 = \qty{1.51e7}{m^3}$. (ب) $0.30 \times 1.512 \times 10^7 =
\qty{4.54e6}{m^3}$ و$0.40 \times 1.512 \times 10^7 = \qty{6.05e6}{m^3}$.
\end{solution}

\begin{solution}{exo:b2:energy-geology:4}
(أ) $V = 10^7 \times 200 = \qty{2e9}{m^3}$ و$q_R = 2.5 \times 10^6 \times 2 \times 10^9 \times 60
= \qty{3.0e17}{J}$. (ب) دون \qty{90}{\celsius}: \omterm{def:b2:energy-geology:geothermal}{مورد منخفض الإنثالبي}، للاستعمال المباشر
للحرارة. (ج) $P = 0.25 \times 3.0 \times 10^{17}/(30 \times \qty{3.156e7}{s}) = \qty{79.2}{MW}$ من الحرارة.
\end{solution}

\begin{solution}{exo:b2:energy-geology:5}
(أ) $k_2/k_1 = \exp[(E/R)(1/T_1 - 1/T_2)] = 2$ مع $1/373.15 - 1/383.15 = \qty{6.994e-5}{K^{-1}}$
يعطي $E = 8.314 \times \ln 2/(6.994 \times 10^{-5}) = \qty{82.4}{kJ/mol}$. (ب) $\exp(26800 \times
6.994 \times 10^{-5}) = 6.52$. (ج) قاعدة التضاعف، وطاقتها الفعلية أدنى، تعطي وزنًا
أكبر نسبيًا للأزمنة الطويلة التي قضاها الصخر في درجة حرارة منخفضة؛ ولكيروجين من النوع~I ذي طاقة
تنشيط عالية، يعلن مؤشّر TTI التوليد مبكّرًا أكثر ممّا ينبغي.
\end{solution}

\begin{solution}{exo:b2:energy-geology:6}
(أ) وضع $x = 1/2$ في التقريب يعطي $(ART^2/\beta E)\,e^{-E/RT} = \ln 2$؛ وبأخذ اللوغاريتمات،
$E/RT = \ln[AT^2/(\beta(E/R)\ln 2)]$، ومنها تنتج الصيغة المطلوبة. وبالبدء من أيّ تخمين قرب
\qty{400}{K} والتكرار مع $\beta$ بالكلفن في الثانية: عند \qty{1}{\celsius\per\Myr}
($\beta = \qty{3.169e-14}{K/s}$)، $T = \qty{409.6}{K}$، أي \qty{136.4}{\celsius}؛ وعند \qty{10}{K} في
الدقيقة ($\beta = \qty{0.1667}{K/s}$، أي أسرع $\num{5.26e12}$ مرّة)، $T = \qty{719.0}{K}$، أي
\qty{445.9}{\celsius}. (ب) الفرق \qty{309.5}{K}؛ والتقدير $6.259 \times
\ln(5.26 \times 10^{12}) = 6.259 \times 29.29 = \qty{183}{K}$ صغير جدًّا، لأنّ $RT^2/E$
نفسه يكبر مع $T$ (فيبلغ \qty{19.3}{K} عند \qty{719}{K}): والإزاحة الخطّية لا تصحّ إلا للتغيّرات
الصغيرة في المعدّل.
\end{solution}

\begin{solution}{exo:b2:energy-geology:7}
$\lambda_5 - \lambda_8 = \qty{8.29725e-10}{yr^{-1}}$ و$r_0 = 0.7204/99.2742 = 0.0072567$.
(أ) قبل \qty{2.0}{\Ga}، $r = 0.0072567\,e^{1.6595} = 0.03814$، والحصّة $0.03814/1.03814 =
\qty{3.67}{\percent}$؛ وقبل \qty{1.8}{\Ga}، $r = 0.03231$، والحصّة \qty{3.13}{\percent}. (ب) $r =
3/97 = 0.030928$ و$t = \ln(0.030928/0.0072567)/(8.29725 \times 10^{-10}) = \qty{1.75}{\Ga}$.
(ج) تحلّل كلّ من الخام المستنفَد واليورانيوم العادي منذ ذلك الحين بالعوامل نفسها، فنسبة
حصّتيهما من $^{235}$U اليوم هي نفسها بُعيد التفاعلات: كان الخام حينئذ يحوي نصف
الحصّة العادية، فقد انشطر نحو نصف ما فيه من $^{235}$U.
\end{solution}

\begin{solution}{exo:b2:energy-geology:8}
(أ) $10^9/1.872 = \qty{5.34e8}{m^3}$؛ و$10^9/621.4 = \qty{1.61e6}{m^3}$؛ و$10^9/686.4 =
\qty{1.46e6}{m^3}$. (ب) $7 \times 10^9/583 = \qty{1.20e7}{m^3}$ من \ce{CO2}، تملأ
$0.35 \times 0.30 = 0.105$ من حجم الصخر: أي $1.20 \times 10^7/0.105 = \qty{1.14e8}{m^3}$ من الصخر،
بمتوسّط سمك قدره $1.143 \times 10^8/(5 \times 10^6) = \qty{22.9}{m}$.
\end{solution}

\begin{solution}{exo:b2:energy-geology:9}
(أ) $V = \qty{9e9}{m^3}$ و$q_R = 2.6 \times 10^6 \times 9 \times 10^9 \times 235 =
\qty{5.50e18}{J}$. (ب) على مدى $30 \times \qty{3.156e7}{s} = \qty{9.47e8}{s}$: تعطي $R_g = 0.1$
\qty{581}{MW} وتعطي $0.25$ \qty{1450}{MW} من الحرارة. (ج) $6 \times (10 \text{ إلى } 20) = 60$ إلى
\qty{120}{MW} من الكهرباء، أي $60/581 = 0.10$ إلى $120/581 = 0.21$ من القدرة الحرارية: فمعظم
الحرارة يُطرح، كما في أيّ محطّة طاقة تعمل بفرق متواضع في درجة الحرارة.
\end{solution}

\begin{solution}{exo:b2:energy-geology:10}
(أ) المقاومة $0.6 \times (80 - 39) = \qty{24.6}{MPa} > \qty{22}{MPa}$: الصدع مستقرّ. ويتطلّب الانهيار
$0.6\,\Delta p = 24.6 - 22$، أي $\Delta p = \qty{4.33}{MPa} = \qty{43}{bar}$، وهو من رتبة
\qty{32}{bar} التي حدّت زلازل دنفر. (ب) $(24.6 - 24)/0.6 = \qty{1.0}{MPa} =
\qty{10}{bar}$: فالصدع القريب أصلًا من الانهيار لا يحتاج إلا إلى ارتفاع صغير. (ج) يرفع المائع المحقون
الضغط في منطقة تتّسع بالانتشار عبر الصخر النفوذ ومناطق الصدوع، فتبلغ كيلومترات
على مدى سنوات؛ وحيثما صادفت صدعًا قريبًا أصلًا من الانهيار، يكفي ارتفاع متواضع.
\end{solution}

\begin{solution}{exo:b2:energy-geology:11}
(أ) اكتب $x = e^{-u}$ مع $u = (ART^2/\beta E)\,e^{-E/RT}$. فنسبة تحوّل قدرها
\qty{10}{\percent} تعني $x = 0.9$، أي $u = \ln(10/9)$؛ و\qty{90}{\percent} تعني $u = \ln 10$. وبما أنّ
$\mathrm{d}\ln u/\mathrm{d}T = E/RT^2 + 2/T \approx E/RT^2$، يكون $\Delta(\ln u) \approx (E/RT^2)\,\Delta T$،
ومنه الصيغة، مع $\ln\ln 10 - \ln\ln(10/9) = 0.834 + 2.250 = 3.084$. (ب) مع $T =
\qty{409.6}{K}$، $RT^2/E = \qty{6.26}{K}$ و$\Delta T = \qty{19.3}{K}$؛ ويعطي المنحنى الدقيق ذو الطاقة
الوحيدة \qty{18.6}{K}، من 125.3 إلى \qty{143.8}{\celsius}. (ج) يرتفع المنحنى المتّصل من
\qty{10}{\percent} قرب \qty{122}{\celsius} إلى \qty{90}{\percent} قرب \qty{148}{\celsius}، على نحو
\qty{26}{K}: فتوزّع طاقات التنشيط يوسّع فترة التحوّل.
\end{solution}

\begin{solution}{exo:b2:energy-geology:12}
(أ) $2 \times 10^{10}/20 = \qty{1.0e9}{m^3}$ (\qty{1.0}{km^3}) و$2 \times 10^{10}/60 =
\qty{3.3e8}{m^3}$. (ب) $2 \times 10^{10}/621.4 = \qty{3.22e7}{m^3}$. (ج) $3.22 \times 10^7/0.02 =
\qty{1.61e9}{m^3}$ من المسام، وهذا مع $50 \times 0.30 = \qty{15}{m}$ من فراغ المسام لكلّ متر
مربّع يتطلّب $1.61 \times 10^9/15 = \qty{1.07e8}{m^2}$، أي نحو \qty{107}{km^2}.
\end{solution}

\begin{solution}{pb:b2:energy-geology:1}
\noindent\textbf{الجزء الأول.}
\textbf{1.} عند \qty{70}{\Ma}: $50 \times 60 = \qty{3000}{m}$، و$15 + 25 \times 3.0 =
\qty{90}{\celsius}$. واليوم: $3000 + 70 \times 30 = \qty{5100}{m}$، و\qty{142.5}{\celsius}.
\textbf{2.} $60 \times 0.025 = \qty{1.5}{\celsius\per\Myr}$، ثم $30 \times 0.025 =
\qty{0.75}{\celsius\per\Myr}$.
\textbf{3.} 15--\qty{20}{\celsius}: \qty{3.33}{\Myr}؛ وكلّ فترة من 20 إلى \qty{90}{\celsius}:
$10/1.5 = \qty{6.67}{\Myr}$؛ وكلّ فترة من 90 إلى \qty{140}{\celsius}: $10/0.75 = \qty{13.3}{\Myr}$؛
و140--\qty{142.5}{\celsius}: \qty{3.33}{\Myr}.
\textbf{4.} قيم الفترات $\Delta t_n 2^n$، مع المجاميع التراكمية: 15--20 ($n = -9$) 0.0065؛ و20--30
0.026 (0.033)؛ و30--40 0.052 (0.085)؛ و40--50 0.104 (0.189)؛ و50--60 0.208 (0.397)؛ و60--70 0.417
(0.814)؛ و70--80 0.833 (1.65)؛ و80--90 1.67 (3.31)؛ و90--100 6.67 (9.98)؛ و100--110 13.3 (23.3)؛
و110--120 26.7 (50.0)؛ و120--130 53.3 (103)؛ و130--140 107 (210)؛ و140--142.5 ($n = 4$) 53.3: فمؤشّر TTI
اليوم 263.
\textbf{5.} تُتجاوز القيمة 15 في الفترة 100--\qty{110}{\celsius}، عند $100 + 10 \times
(15 - 9.98)/13.33 = \qty{103.8}{\celsius}$، التي بُلغت بعد $(103.8 - 90)/0.75 = \qty{18.4}{\Myr}$ من
\qty{70}{\Ma}: أي عند \qty{51.6}{\Ma} وعلى عمق $(103.8 - 15)/0.025 = \qty{3550}{m}$.
\textbf{6.} 75: \qty{124.7}{\celsius}، و\qty{23.7}{\Ma}، و\qty{4390}{m}. و160: \qty{135.3}{\celsius}،
و\qty{9.6}{\Ma}، و\qty{4810}{m}.
\textbf{7.} لا: فعند 263 يكون قد تجاوز نهاية توليد النفط (160)، منذ نحو \qty{10}{\Ma}.

\noindent\textbf{الجزء الثاني.}
\textbf{8.} $k = 9.21 \times 10^{13}\,e^{-26800/363.15} = \qty{8.20e-19}{s^{-1}}$. كان الصخر
دون \qty{90}{\celsius} طوال \qty{50}{\Myr} التي سبقت \qty{70}{\Ma}، فالكسر المتحوّل
أقلّ من $k\,t = 8.20 \times 10^{-19} \times 50 \times 3.156 \times 10^{13} = \num{1.29e-3}$.
\textbf{9.} $\beta = 0.75/(3.156 \times 10^{13}) = \qty{2.376e-14}{K/s}$، و$T = \qty{415.65}{K}$:
$ART^2/(\beta E) = \num{2.498e28}$ و$e^{-E/RT} = \num{9.95e-29}$، فيكون $u = 2.486$ و$x = 0.0833$
ونسبة التحوّل 0.917.
\textbf{10.} \qty{134.7}{\celsius}.
\textbf{11.} $70 - (134.7 - 90)/0.75 = \qty{10.4}{\Ma}$.
\textbf{12.} يضع مؤشّر TTI ذروة التوليد عند \qty{23.7}{\Ma}، والحركية بعد ذلك بمقدار \qty{13}{\Myr}،
لأنّ قاعدة التضاعف تقابل $E \approx \qty{82}{kJ/mol}$ وتعطي درجات الحرارة المنخفضة وزنًا
أكبر ممّا ينبغي لكيروجين طاقته \qty{223}{kJ/mol} (التمرين~\ref{exo:b2:energy-geology:5}).
\textbf{13.} في وقت أبكر: فالطاقة الأدنى تخفض درجة حرارة نصف التحوّل.

\noindent\textbf{الجزء الثالث.}
\textbf{14.} الصخر: $4 \times 10^8 \times 50 \times 2400 = \qty{4.8e13}{kg}$؛ والإمكانية $0.030
\times 4.8 \times 10^{13} = \qty{1.44e12}{kg}$؛ والمولَّد $0.92 \times 1.44 \times 10^{12} =
\qty{1.32e12}{kg}$.
\textbf{15.} $1.3248 \times 10^{12}/850 = \qty{1.56e9}{m^3}$.
\textbf{16.} $1.2 \times 10^7 \times 40 \times 0.22 \times 0.75/1.30 = \qty{6.09e7}{m^3}$.
\textbf{17.} $6.09 \times 10^7/1.56 \times 10^9 = 0.039$، أي \qty{3.9}{\percent}: فمعظم النفط المولَّد
يبقى في الصخر المصدر، أو يضيع على طول مسارات الهجرة، أو يفلت إلى السطح.
\textbf{18.} $0.35 \times 6.09 \times 10^7 = \qty{2.13e7}{m^3}$.
\textbf{19.} يبقى نحو ثلثي النفط في المكمن بعد الاستخلاص الأوّلي والثانوي؛
ويمكن أن يستخلص حقن \ce{CO2} من 10 إلى \qty{15}{\percent} إضافية من النفط في المكمن.

\noindent\textbf{الجزء الرابع.}
\textbf{20.} $15 + 25 \times 2 = \qty{65}{\celsius}$ و$0.101 + 1000 \times 9.807 \times 2000/10^6
= \qty{19.7}{MPa}$: كلاهما فوق النقطة الحرجة (\qty{31.0}{\celsius}، \qty{7.38}{MPa})، فيكون
\ce{CO2} فوق حرج، بكثافة نحو \qty{686}{kg.m^{-3}}.
\textbf{21.} $2.13 \times 10^7 \times 1.30 = \qty{2.77e7}{m^3}$.
\textbf{22.} حجم المسام $1.2 \times 10^7 \times 40 \times 0.22 = \qty{1.056e8}{m^3}$، و$E =
2.77 \times 10^7/1.056 \times 10^8 = 0.262$، وهذا أعلى بكثير من 1 إلى \qty{4}{\percent} في التكوينات
الملحية: فالنفط المنتَج أفرغ حيّزًا معروفًا ومُحكم الإغلاق.
\textbf{23.} $6 \times 10^8 \times 80 \times 0.20 \times 691 \times 0.02 = \qty{1.33e11}{kg}$،
أي \qty{133}{Mt}.
\textbf{24.} بمعدّل \qty{1}{Mt} في السنة، نحو 133 سنة.
\textbf{25.} $M = \rho\,N_p\,B_o = 686 \times 2.77 \times 10^7 = \qty{1.90e10}{kg}$: يستطيع المكمن
المستنفَد أن يخزّن \textbf{\qty[text-series-to-math=true]{19.0}{Mt}} من \ce{CO2}.
\end{solution}
